\section{Crossings and the scheduler}\label{sec:framework}

Fix a linear coordinate \(x\), which we call time. We want to force many
changes of direction on a path from one time level to another. The basic
observation is that a point on a face of dimension at most \(J\) admits
at most \(J\) independent two-sided directions. The recursive construction
uses this observation twice. It first constrains a path in an outer
polyhedron, then reads a path in an inner polyhedron from slacks of the
outer inequalities. Only some inner coordinates can be varied independently
through these slacks. We record which sets of coordinates suffice for the
next crossing.

\subsection{Blocking inequalities and crossings}

A \emph{blocking system} \(Q\) consists of finitely many inequalities
\begin{equation}\label{eq:blocking-row}
             c\cdot z\geq \ell(x),\qquad c\geq0,\quad c\ne0,
\end{equation}
with affine right-hand sides. It contains an individual lower bound for
each coordinate of \(z\). A convex piecewise affine right-hand side stands
for all the inequalities given by its affine pieces. Rows are indexed,
so that equal inequalities with different indices retain their identities.
The \emph{support} of a row is \(\{i:c_i>0\}\), and its \emph{slack} at
\((x,z)\) is \(c\cdot z-\ell(x)\).

Write \(Q(x)\) for the slice at time \(x\). For a set \(I\) of coordinates,
restrict the system by keeping only coordinates in \(I\) and rows supported
on \(I\). The restricted slice is exactly the projection of \(Q(x)\) onto
\(I\). Indeed, a feasible vector on \(I\) extends to \(Q(x)\) by taking
every omitted coordinate equal to a sufficiently large number. Each
discarded row has a positive coefficient on at least one such coordinate,
and there are only finitely many rows. The same argument gives a strictly
feasible point in every restricted slice. Its individual lower bounds
exclude lines.

A \emph{crossing} of \([a,b]\) is a continuous polygonal path in \(Q\)
whose time coordinate starts at \(a\), ends at \(b\), and stays in
\([a,b]\). It may reverse direction in time. A \emph{segmentation} is a
finite partition of the path's parameter interval on each part of which
the whole path is affine. Constant pieces are allowed. A \emph{break} is
an interior partition point. On a subpath we count only breaks strictly
inside its parameter interval, so that neighboring traversals cannot
charge a common endpoint twice.

\subsection{Which coordinates suffice?}

Consider a coordinate \(i\) whose role in the crossing argument can also be
fulfilled by two other coordinates \(j,k\) together. The usable sets
are those containing \(i\) or containing both \(j\) and \(k\). More
generally a coordinate may have several alternatives, each with further
alternatives of its own.

A \emph{signature} \(T\) is a finite rooted tree. Its height is the largest
number of edges on a root-to-leaf path. Each edge \(\gamma\) out of the
root is an \emph{option}, with signature \(T_\gamma\) given by the subtree
rooted at its child. A \emph{bank} is an occurrence of a signature. A
\emph{program} on a bank of signature \(T\) makes a finite list of
\emph{direct requests}. For each such request \(i\) and each option
\(\gamma\), it specifies a program on a fresh bank of signature
\(T_\gamma\), the \emph{upgrade} of \(i\) along \(\gamma\). Lists may be
empty. Each request has its own identity and indexes one coordinate.

Call a set \(I\) of requests \emph{viable} if for every direct request
\(i\), either \(i\in I\) or some upgrade of \(i\) is itself viable for
\(I\). At a bank with no options all direct requests must belong to \(I\).
An empty program is viable for every \(I\). In the example above \(i\) has a
single upgrade, the terminal list \(j,k\). Adding coordinates preserves
viability. We may combine programs at one bank by concatenating disjoint
direct lists and keeping each request's upgrades. The combined program
is viable precisely when each component program is viable.

One useful program makes one direct request at the root and repeats this
prescription in every upgrade, down to the leaves of \(T\). It has one
request per node. A set \(I\) is viable precisely when it contains at
least one of these requests. Thus \(J+1\) disjoint copies at the same
bank require \(J+1\) distinct usable coordinates. If these coordinates
occur only in individual lower bounds and all those bounds are slack,
their \(J+1\) coordinate directions are independent two-sided motions.
This is how the construction in Section~\ref{sec:recursion} forces an
equality on a low-dimensional face.

\begin{definition}[Admissibility]\label{def:admissible}
Let \(J\ge0\) be an integer and let \(Q\) be a blocking system on the
requests of a program. At a feasible point \((x,z)\), form the matrix
\(C\) whose rows are the coordinate normals \(c\) of the tight
inequalities. The point is \emph{admissible} if there exists a
set \(I\) of \emph{usable} coordinates such that
\begin{enumerate}
\item \(I\) is viable for the program,
\item every tight row is supported on \(I\), and
\item \(\dim\ker(C|_I)\le J\), where \(C|_I\) retains the columns in \(I\).
\end{enumerate}
A crossing is admissible if each point is admissible, with \(I\)
allowed to vary from point to point.
\end{definition}

At fixed \(x\) a direction is feasible for small moves in both signs
precisely when it annihilates every tight coordinate normal. The
two-sided directions therefore form the translation space of the
smallest face of \(Q(x)\) containing \(z\). A point of a slice face of
dimension at most \(J\) meets Definition~\ref{def:admissible} with \(I\)
equal to the full coordinate set. For the inner path we can instead use
a smaller viable set. All tight rows are still tested in the full system.
The nullity bound then counts the two-sided directions supported on \(I\).

To count the cost of substituting slacks, also consider the forest whose
vertices are \emph{requests}. The children of a request are all the
direct requests in all of its upgrades. The \emph{span} of a set of
requests is the cardinality of its union with all its ancestors in this
forest. The span of a row is that of its support. When several supports
have common ancestors, their union counts those ancestors only once.

\subsection{Forcing equality at selected times}

The \emph{scheduler} is a blocking system for a prescribed list of
time intervals. Its input includes a program \(\pi\), called the
\emph{client}. For each index \(P\in\{1,\ldots,B\}\), called a
\emph{bucket}, the scheduler supplies a group of distinguished
\emph{red rows} for each client request \(i\). A tight row in the group
\((P,i)\) certifies that \(i\) is usable for bucket \(P\).

Each interval has two endpoints, its \emph{anchors}, and one interior
\emph{probe} per bucket. At either anchor, only a few buckets can have any
nearly tight red row. If the number of buckets exceeds twice this
per-anchor bound, then for any two feasible points at the anchors there
is a bucket with every red row slack at both points. Near its probe, the
tight red rows at a vertex certify a viable set of client requests. This
probe conclusion survives restriction to a viable set of scheduler
coordinates and small perturbations of the right-hand sides.

\begin{theorem}[Scheduler]\label{thm:scheduler}
For every integer $r\geq100$ there are arbitrarily large integers $L$
with an associated label set of size $V$ satisfying $10V\leq L^{3/r}$.
Fix such parameters, take an integer $h\geq1$, and put $N=L^h$. For each signature
$T$ of height $H$ and maximum outdegree $G$ there is a signature $T^+$ of
height at most $h+2H$ and maximum outdegree at most $10V(G+1)$.

Let $\pi$ be any client program on $T^+$ with $m\geq1$ requests in total,
and let $B\geq1$ be an integer. Prescribe $N$ pairwise disjoint closed intervals
(\emph{episodes}) in increasing order. In each interval designate its two
endpoints as \emph{anchors} and choose $B$ interior times
(\emph{probes}), one for each bucket in order. There is a blocking system
on a program on $T$ with the following properties.
\begin{enumerate}
\item The red rows are partitioned into groups indexed by $(P,i)$, where
$P$ is a bucket and $i$ a client request. Each group contains between one
and $2V$ rows. Deleting all red rows leaves an individual lower bound on
every scheduler coordinate.
\item Each probe for $P$ lies in the interior of a closed \emph{probe band}.
The bands are pairwise disjoint and lie inside their respective episodes.
Fix any time in the band and restrict the scheduler to any
coordinate set $I$. The coordinate normals tight at any feasible point of
the restricted slice are independent on $I$. If $I$ is viable for the
scheduler, then at every vertex of that slice the set
\[
   \{i:\text{a retained red row in group }(P,i)\text{ is tight}\}
\]
is viable for $\pi$. The bands can be chosen so that both assertions hold
after arbitrary changes of the row right-hand sides in the slice by less
than a common positive tolerance.
\item At each anchor there is a tolerance $\xi>0$ such that whenever
every scheduler row has slack at least $-\xi$, at most $2Vm$ buckets
contain a red row of slack at most $\xi$. Here the slack of a row
in~\eqref{eq:blocking-row} is $c\cdot z-\ell(x)$.
\item There are at most $10VBm$ scheduler coordinates and
\[
               O_{r,L,V}\!\left(NBm\,[1+(Bm)^2]\right)
\]
rows. Each row has span at most $C(h+H)+O_{r,L,V}(1)$ for an absolute
constant $C$.
\item Fix a bucket and any set of client requests of span $s$. Choose one
red row in the designated bucket for each request in the set. The union
of their supports has span at most $5s$ in the scheduler program.
\end{enumerate}
The signature $T^+$ is independent of the client program, the number of
buckets, and the prescribed times.
\end{theorem}

The two slice assertions in Part~2 have different roles. At a point on a
positive-dimensional face in the probe band, we can complete its tight
equalities to those of a vertex in an appropriate restricted slice. The
vertex supplies viable client requests. Independence counts the additional
equalities needed for this completion. At the anchors the guarantee comes
from Part~3 and applies even at a slightly infeasible point of the full
system. Part~5 controls the number of rows generated when a minimum of
red slacks is substituted for each coordinate of an inner inequality.
Sections~\ref{sec:layers}--\ref{sec:compiler} construct the scheduler
and prove Theorem~\ref{thm:scheduler}.
